% Generated draft: requires mathematical and source review.
\documentclass[11pt]{article}
\usepackage[margin=25mm]{geometry}
\usepackage{fontspec}
\setmainfont{Latin Modern Roman}
\usepackage{amsmath,amssymb,mathtools,unicode-math}
\setmathfont{Latin Modern Math}
\usepackage{enumitem,needspace,xurl,hyperref}
\setlength{\parindent}{0pt}
\setlength{\parskip}{5pt}
\setlength{\emergencystretch}{3em}
\newcommand{\N}{\mathbb{N}}
\newcommand{\Z}{\mathbb{Z}}
\newcommand{\Q}{\mathbb{Q}}
\newcommand{\R}{\mathbb{R}}
\newcommand{\C}{\mathbb{C}}
\newcommand{\F}{\mathbb{F}}
\newcommand{\A}{\mathbb{A}}
\newcommand{\PP}{\mathbb{P}}
\newcommand{\E}{\mathbb{E}}
\newcommand{\eps}{\varepsilon}
\newcommand{\dd}{\,\mathrm{d}}
\begin{document}
\section*{UnsolvedMath ID: 700337}
\textbf{Source problem code:} catalogue entry
\subsection{Definitions and mathematical statement}
For a group \(G\), write
\[
[G,G]=\langle [x,y]:x,y\in G\rangle,
\qquad [x,y]=x^{-1}y^{-1}xy.
\]
This is the derived subgroup of \(G\). A group is perfect if \(G=[G,G]\). A group is locally nilpotent if every finitely generated subgroup is nilpotent. Here, exponent \(p\) means that every element satisfies \(g^p=1\) and at least one element has order \(p\). The alternative convention allowing the trivial group does not change either existence question below.

\textbf{Part (a).} Is the following assertion true?
\[
\begin{gathered}
\text{For every prime }p>3\text{ and every finite group }G
\text{ of exponent }p,\\
\text{there exist a finite group }H\text{ of exponent }p
\text{ and an injective homomorphism }\iota:G\longrightarrow[H,H].
\end{gathered}
\]
The supplied registry places this part among the remaining parts of the problem; this is a supplied status assessment, not an independently verified claim of openness.

\textbf{Part (b), archived as solved in the supplied record.} Does there exist a prime \(p\) and a nontrivial locally nilpotent group \(G\) of exponent \(p\) such that
\[
G=[G,G]?
\]
The question quantifies existentially over the prime; it does not ask for such a group for every prime. Nontriviality is explicit in the accompanying archive account, although it is omitted from the wording of the question itself. It is therefore included here as the intended qualification. The source adds that such a group has no maximal subgroups, where a maximal subgroup means a proper subgroup contained in no larger proper subgroup.

\textbf{Supplied historical background.} The archive reports an affirmative answer to part (b). Its stated construction uses a non-soluble variety \(\mathfrak V\), a class of groups specified by identities with no common finite bound on derived length. Let \(F\) be the free group in \(\mathfrak V\) on generators \(x_1,x_2,\ldots\), and let \(\varphi:F\to F\) be the homomorphism determined by
\[
\varphi(x_i)=[x_{2i-1},x_{2i}]\qquad(i\geq1).
\]
The archive states that
\[
D=\varinjlim\bigl(F\xrightarrow{\varphi}F\xrightarrow{\varphi}F\xrightarrow{\varphi}\cdots\bigr)
\]
is nontrivial and satisfies \(D=[D,D]\). It further reports that the Kostrikin variety of locally nilpotent groups of prime exponent \(p\geq5\) is non-soluble, citing Razmyslov and also Kolmakov, and credits a letter of Khukhro and L'vov with the account. It cites Vaughan-Lee and Wiegold for another example answering the question. These are preserved source reports, not new proofs or independently verified literature conclusions; the record does not provide a separate definition of the named Kostrikin variety.
\subsection{Short English statement}
Part (a) asks whether, for every prime \(p>3\), each finite group of exponent \(p\) embeds into the derived subgroup of another finite group of exponent \(p\). Part (b) asks for a nontrivial locally nilpotent perfect group of some prime exponent. The supplied record labels the whole problem partially solved and archives an affirmative answer to part (b); these labels and the historical claims require human source review.
\subsection{Sources}
\begin{itemize}
\item Supplied local UnsolvedMath registry record, permanent ID \(700337\), exact source problem code \texttt{catalogue entry}, titled \emph{Kourovka Notebook Problem 5.56}. The record attributes the problem to J. Wiegold, proposed in 1976, and labels it \texttt{partially\_solved}.
\item \emph{The Kourovka Notebook (Unsolved Problems in Group Theory)}, 5th issue (1976), Problem 5.56; supplied background describes an update in the 21st issue (2026), arXiv identifier \texttt{1401.0300v48}. The supplied collection description names E. I. Khukhro and V. D. Mazurov as editors. Part (b) is supplied as an Archive of Solved Problems entry, while part (a) remains in the main section according to the registry review.
\item Yu. P. Razmyslov, \emph{Algebra and Logic}, volume 10 (1971), pages 21--29. The archive attributes the non-solubility assertion concerning the Kostrikin variety to this reference. No article title or link is supplied.
\item E. I. Khukhro and I. V. L'vov, letter dated June 19, 1976, cited in the supplied archive account. No further publication details or link are supplied.
\item Yu. A. Kolmakov, \emph{Math. Notes}, volume 35, nos. 5--6 (1984), pages 389--391. The archive reports that the same non-solubility assertion was also proved there. No article title or link is supplied.
\item M. R. Vaughan-Lee and J. Wiegold, \emph{Bull. London Math. Soc.}, volume 13, no. 1 (1981), pages 45--46. The archive cites this for another example answering part (b). No article title or link is supplied.
\item The registry status review is dated \texttt{2026-10-16}, with evidence labeled \texttt{published}. This is supplied dated curation, not an independently verified literature assessment. No separate \texttt{literature\_checked\_at} field or \texttt{literature\_sources} links are supplied.
\end{itemize}
\end{document}
